\section{Virtual-sequence approach -  Mishchenko
et al. (2022)}

\subsection{Setup}

\begin{itemize}
    \item Assume $B=1$ while keeping $G\geq 2$ $\implies$ at every learner iteration, exactly one completed
group is consumed
    \item $M$ rollout workers, each generating one group at a time. when a group finishes, the learner immediately applies its GMC update and the worker starts generating a new group of responses using the updated policy.
    \item Every dispatched group eventually returns and is consumed once
    \item Completion times depend on generated responses
\end{itemize}


\subsubsection{Correspondence with Mishchenko et. al.}

\begin{tabular}[t]{|l|l|}
\hline Mishchenko et al. & Our setting \\
\hline Parameter vector $x_k$ & Policy parameters $\theta_n$ at learner iteration $n$ \\
\hline Worker $m_k$ & Rollout worker $m_k$ \\
\hline Random sample $\xi_k^{m_k}$ & Prompt and responses $\left(X_k, Y_k^1, \ldots, Y_k^G\right)$ \\
\hline Gradient computed at $x_k$ & Group generated under $\pi_{\theta_{d(k)}}$ \\
\hline Stochastic gradient $\nabla f\left(x_k ; \xi_k^{m_k}\right)$ & Generation-time estimator $V_k$ \\
\hline
\end{tabular}

Structural differences:
\begin{itemize}
    \item  Mishchenko's worker computes a gradient, whereas our worker generates responses, and actual GMC gradient is evaluated later.
    \item In their Algorithm 1, a worker starts computing
    $
    \nabla f\left(x_k ; \xi_k^{m_k}\right)
    $
    using the parameters $x_k$. If this gradient finishes at iteration $r(k)$, the learner appli that same gradient:

    $$
    x_{r(k)+1}=x_{r(k)}-\eta \nabla f\left(x_k ; \xi_k^{m_k}\right) .
    $$

    The gradient is stale, but its value is not recomputed at the current parameters. In GMC-GRPO, the responses are generated at dispatch, but the actual gradient estimator is evaluated when the group returns.
    \item Unlike in Mishchenko et al., the applied update is not identical to the generation-time estimator (see below).
    \end{itemize}


\subsubsection{Dispatch-time sampling}

Suppose group $k$ is dispatched at iteration $d(k)$ and consumed at iteration $r(k)$.

Let $\mathcal G_k$ denote the information available immediately
before this group is sampled. At dispatch, we assume
\begin{equation*}
X_k\sim\mathcal D,
\qquad
Y_k^1,\ldots,Y_k^G
\mid X_k,\mathcal G_k
\overset{\mathrm{i.i.d.}}{\sim}
\pi_{\theta_{d(k)}}(\cdot\mid X_k).
\end{equation*}

We allow the completion time to depend on the generated
responses. In particular, no independence between
response lengths and delays is assumed.

Define the generation-time estimator
\begin{equation*}
V_k=
\frac{1}{TG}\sum_{i=1}^G
\nabla\log\pi_{\theta_{d(k)}}(Y_k^i\mid X_k)
\widehat A_k^i.
\end{equation*}

Then
\begin{equation*}
\mathbb E[V_k\mid\mathcal G_k]
=\frac1T\nabla J(\theta_{d(k)}).
\end{equation*}

This is the analogue of the unbiased stochastic gradient
$\nabla f(x_k;\xi_k^{m_k})$ used in Mishchenko et al.
The expectation is taken at dispatch, before conditioning
on the group's completion time, so no independence assumption needed.

\paragraph{Difference between $\mathcal F_n^-$ and $\mathcal G_k$ - why conditional unbiasedness now holds}
\begin{itemize}
    \item $\mathcal F_n^-$ contains the information available before the next
completed group is consumed,
    \item $\mathcal G_k$ contains the
information available before a new group is generated
\end{itemize}
The group arriving at iteration $n$ has already been generated, and (the order of) its arrival may depend on response
length, so
\begin{equation*}
\mathbb E[g_n\mid\mathcal F_n^-]
\neq \frac1T\nabla J(\theta_n).
\end{equation*}


At dispatch, the group is sampled from
$\pi_{\theta_{d(k)}}$, before its completion time
is observed. The completion time may depend on
the generated responses. Nevertheless,
\[
\mathbb E[V_k\mid\mathcal G_k]
=\frac1T\nabla J(\theta_{d(k)}),
\]
since we do not condition on the completion time.

\subsubsection{Return-time GMC estimator}

Suppose group $k$ is dispatched at iteration $d(k)$ and consumed at iteration $r(k)$.

$$
\begin{gathered}
\rho_k^i=\frac{\pi_{\theta_{r(k)}}\left(Y_k^i \mid X_k\right)}{\pi_{\theta_{d(k)}}\left(Y_k^i \mid X_k\right)}, \quad \overline{\rho_k}=\frac{1}{G} \sum_{i=1}^G \rho_k^i, \\
\hat{\rho}_k^i=\frac{\rho_k^i}{\max \left\{1, \bar{\rho}_k /(1+\epsilon)\right\}} .
\end{gathered}
$$

The learner applies

$$
U_k=\frac{1}{T G} \sum_{i=1}^G \hat{\rho}_k^i \nabla \log \pi_{\theta_{r(k)}}\left(Y_k^i \mid X_k\right) \widehat{A}_k^i .
$$

(same as before, but indexed the group by dispatch instead of consumption iteration)

Define

$$
e_k:=U_k-V_k .
$$

This is the effect of evaluating the group update at the return-time policy rather than at its generation-time policy.
Then

$$
U_k=V_k+e_k .
$$

\subsubsection{Consumed, dispatched and oustanding groups}
Consider the moment immediately after the learner has completed $n$ updates and dispatched the next group.

Define

$$
\begin{aligned}
\mathcal{D}_n & =\{\text { groups dispatched so far }\}, \\
\mathcal{C}_n & =\{\text { groups consumed so far }\}, \\
\mathcal{P}_n & =\mathcal{D}_n \backslash \mathcal{C}_n .
\end{aligned}
$$


Here $\mathcal{P}_n$ is the set of outstanding groups.
Since every worker generates one group at a time,

$$
\left|\mathcal{P}_n\right|=M .
$$

\subsection{Virtual iterates}

The actual learner has applied only the completed groups:

$$
\theta_n=\theta_0+\eta \sum_{j \in C_n} U_j .
$$

Define a virtual learner that has already applied the generation-time estimators of all dispatched groups (also those still generating):

$$
\tilde{\theta}_n=\theta_0+\eta \sum_{j \in \mathcal{D}_n} V_j .
$$
The real learner cannot apply a group's estimator before its responses are available and applies the return-time, not generation-time estimator.

\begin{lemma}[Analogue of Mishchenko et al.'s Lemma 1]
    \begin{align*}
    \tilde{\theta}_n-\theta_n & =\eta \sum_{j \in \mathcal{D}_n} V_j-\eta \sum_{j \in C_n} U_j \\
    & =\eta \sum_{j \in \mathcal{D}_n} V_j-\eta \sum_{j \in C_n}\left(V_j+e_j\right) \\
    & =\eta \sum_{j \in \mathcal{P}_n} V_j-\eta \sum_{j \in C_n} e_j .
    \end{align*}
    From our Lemma 2 and $\left|\widehat{A}_j^i\right| \leq 2 R$,
    \begin{equation} \label{eq:gentimeestbound}
            \left\|V_j\right\| \leq K_0, \quad K_0=2 R \sqrt{2 L \log C} .
    \end{equation}


    Since $\left|\mathcal{P}_n\right|=M$,
    $
    \left\|\sum_{j \in \mathcal{P}_n} V_j\right\| \leq M K_0 .
    $
    Consequently,

    $$
    \left\|\tilde{\theta}_n-\theta_n\right\| \leq \eta M K_0+\eta \sum_{j \in C_n}\left\|e_j\right\| .
    $$
\end{lemma}

Main difference: their difference consists of outstanding gradients; our difference consists of outstanding estimators and accumulated $e_j$ errors.


\subsection{Bounding $e_k$}


For a fixed generated group, define
\begin{equation*}
F_k(\theta)
=
\frac{1}{TG}\sum_{i=1}^G
\widehat\rho_k^i(\theta)
\nabla\log\pi_\theta(Y_k^i\mid X_k)
\widehat A_k^i,
\end{equation*}
where $$\rho_k^i(\theta)=\frac{\pi_\theta\left(Y_k^i \mid X_k\right)}{\pi_{\theta_{d(k)}}\left(Y_k^i \mid X_k\right)}.$$

Since all ratios equal one at $\theta_{d(k)}$,
\begin{equation*}
F_k(\theta_{d(k)})=V_k,
\qquad
F_k(\theta_{r(k)})=U_k.
\end{equation*}
Thus
\begin{equation*}
e_k=F_k(\theta_{r(k)})-F_k(\theta_{d(k)}).
\end{equation*}

\paragraph{Assumptions}
\begin{itemize}
    \item For each group $k$, along the line segment
\begin{equation*}
\theta(t)=\theta_{d(k)}
+t(\theta_{r(k)}-\theta_{d(k)}),
\qquad t\in[0,1],
\end{equation*}
the GMC cap switches between active and inactive only
finitely many times.
\item 
We additionally assume that Assumptions 1 and 2
hold along every line segment joining
$\theta_{d(k)}$ and $\theta_{r(k)}$.
\end{itemize}

\begin{lemma}[Error bound] \label{lem:gmc_error}
    $$\|e_k\| \le H_\epsilon\eta K\tau_k$$
    where
    $
    H_\epsilon
    =
    2R(1+\epsilon)(4TL\log C+L),
    $
    $
    K=2R(1+\epsilon)\sqrt{2L\log C},$ and
    $\tau_k=r(k)-d(k).$
\end{lemma}

\begin{proof}

\textbf{Step 1.} Write
\begin{equation*}
S_i(\theta):=\nabla\log\pi_\theta(Y_k^i\mid X_k),
\qquad
S:=T\sqrt{2L\log C}.
\end{equation*}
By Lemma 2 and Assumption 1 (assuming the score bound holds along the parameter segment),
\begin{equation*}
\|S_i(\theta)\|_2\le S,
\qquad
\|S_i(\theta)-S_i(\theta')\|_2
\le TL\|\theta-\theta'\|_2.
\end{equation*}
Since the behavior policy is fixed,
\begin{equation*}
\nabla\rho_k^i(\theta)=\rho_k^i(\theta)S_i(\theta).
\end{equation*}


\textbf{Step 2.} We bound the derivatives of the GMC weights with respect to $\theta$, keeping the behavior policy and generated trajectories fixed.

Recall
\begin{equation*}
\rho_k^i(\theta)
=
\frac{\pi_\theta(Y_k^i\mid X_k)}
{\pi_{\theta_{d(k)}}(Y_k^i\mid X_k)},
\qquad
\bar\rho_k(\theta)=\frac1G\sum_{j=1}^G\rho_k^j(\theta).
\end{equation*}
Since the denominator is fixed,
\begin{equation*}
\nabla_\theta\rho_k^i(\theta)
=\rho_k^i(\theta)S_i(\theta),
\qquad
\nabla_\theta\bar\rho_k(\theta)
=\frac1G\sum_{j=1}^G\rho_k^j(\theta)S_j(\theta).
\end{equation*}

If $\bar\rho_k(\theta)<1+\epsilon$, then
$\hat\rho_k^i=\rho_k^i$, so
\begin{equation*}
\nabla_\theta\hat\rho_k^i
=\hat\rho_k^i S_i
\quad\Rightarrow\quad
\|\nabla_\theta\hat\rho_k^i\|_2
\le S\hat\rho_k^i.
\end{equation*}

If $\bar\rho_k(\theta)>1+\epsilon$, then
\begin{equation*}
\hat\rho_k^i
=(1+\epsilon)\frac{\rho_k^i}{\bar\rho_k}.
\end{equation*}
Applying the quotient rule,
\begin{align*}
\nabla_\theta\hat\rho_k^i
&=(1+\epsilon)
\left(
\frac{\nabla_\theta\rho_k^i}{\bar\rho_k}
-\frac{\rho_k^i\nabla_\theta\bar\rho_k}
{\bar\rho_k^2}
\right)\\
&=\hat\rho_k^i
\left(
S_i-
\frac{\sum_{j=1}^G\rho_k^j S_j}
{\sum_{j=1}^G\rho_k^j}
\right).
\end{align*}

The second term is a weighted average of trajectory scores, whose norm is at most $S$. Hence
\begin{equation*}
\|\nabla_\theta\hat\rho_k^i\|_2
\le \hat\rho_k^i(\|S_i\|_2+S)
\le 2S\hat\rho_k^i.
\end{equation*}

Combining both cases, away from cap-switching points,
\begin{equation}
\|\nabla_\theta\hat\rho_k^i(\theta)\|_2
\le 2S\hat\rho_k^i(\theta).
\label{eq:weight_derivative_bound}
\end{equation}

\textbf{Step 3.} By the product rule, using
$|\widehat A_k^i|\le 2R$, the score bounds and
\eqref{eq:weight_derivative_bound}, wherever the derivative exists,
\begin{align*}
\|DF_k(\theta)\|_{\mathrm{op}}
&\le
\frac{2R}{TG}\sum_{i=1}^G
\left(
\|\nabla\hat\rho_k^i\|_2\|S_i\|_2
+\hat\rho_k^i TL
\right)\\
&\le
\frac{2R}{TG}\sum_{i=1}^G
\hat\rho_k^i(2S^2+TL)\\
&\le
\frac{2R(1+\epsilon)}{T}(2S^2+TL)\\
&=
2R(1+\epsilon)(4TL\log C+L)
=:H_\epsilon,
\end{align*}
where the third inequality uses
$\frac1G\sum_i\hat\rho_k^i\le 1+\epsilon$.

By the finite cap-switching assumption, the segment between
$\theta_{d(k)}$ and $\theta_{r(k)}$ can be partitioned into
finitely many intervals on which the cap has a fixed regime.
Integrating the derivative bound on these intervals and using
continuity at the switching points gives
\begin{equation*}
\|F_k(\theta_{r(k)})-F_k(\theta_{d(k)})\|_2
\le H_\epsilon
\|\theta_{r(k)}-\theta_{d(k)}\|_2.
\end{equation*}
Hence
\begin{equation*}
\|e_k\|_2
\le H_\epsilon
\|\theta_{r(k)}-\theta_{d(k)}\|_2.
\end{equation*}

\textbf{Step 4.} Finally, using the learner update
$\theta_{n+1}=\theta_n+\eta g_n$,
\begin{equation*}
\theta_{r(k)}-\theta_{d(k)}
=
\eta\sum_{n=d(k)}^{r(k)-1}g_n.
\end{equation*}

(Note: $g_n=U_{k(n)}$ where $k(n)$ is group consumed at iteration $n$. Equivalently $g_{r(k)}=U_k$ .)
By adapted Proposition 1,
\begin{equation*}
\|g_n\|_2\le K,
\qquad
K=2R(1+\epsilon)\sqrt{2L\log C}.
\end{equation*}

Then
\begin{equation}
\|\theta_{r(k)}-\theta_{d(k)}\|_2
\le  \eta K(r(k)-d(k)).
\label{eq:parameter_drift}
\end{equation}

Therefore,
\begin{align*}
\|e_k\|_2
&\le H_\epsilon
\|\theta_{r(k)}-\theta_{d(k)}\|_2\\
&\le H_\epsilon\eta K(r(k)-d(k)).
\end{align*}
Writing $\tau_k=r(k)-d(k)$, we obtain
\begin{equation*}
\|e_k\|_2\le H_\epsilon\eta K\tau_k.
\end{equation*}
\end{proof}


\subsection{Bounding sum of delays}

Recall: $d(k)$ and $r(k)$ are the dispatch and consumption iterations of group $k$, and
$$
\tau_k=r(k)-d(k) .
$$
Then
$$
\mathcal{C}_N=\{k: r(k)<N\}
$$
denotes the groups consumed during the first $N$ learner iterations.

\begin{lemma}[Cumulative delay bound] \label{lem:cumulative_delay}
    Suppose at most $M$ groups can be outstanding at any time. Then
    $$
\sum_{k \in C_N} \tau_k \leq M N .
$$
\end{lemma} 

\begin{proof}
For $k\in\mathcal C_N$,
\[
\tau_k
=\sum_{n=0}^{N-1}
\mathbf 1_{\{d(k)\le n<r(k)\}}.
\]
Hence
\begin{align*}
\sum_{k\in\mathcal C_N}\tau_k
&=\sum_{n=0}^{N-1}
\sum_{k\in\mathcal C_N}
\mathbf 1_{\{d(k)\le n<r(k)\}}\\
&\le\sum_{n=0}^{N-1}M
=MN,
\end{align*}
since at most $M$ groups can be outstanding
at any iteration.
\end{proof}

\begin{corollary}
Using the previous bound
$\|e_k\|_2\le H_\epsilon\eta K\tau_k$,
\begin{align*}
\frac1N\sum_{k\in\mathcal C_N}\|e_k\|_2
&\le\frac{H_\epsilon\eta K}{N}
\sum_{k\in\mathcal C_N}\tau_k\\
&\le H_\epsilon\eta KM.
\end{align*}
\end{corollary}

\subsection{Convergence analysis}


\begin{theorem}[Convergence of asynchronous GMC-GRPO]
\label{thm:async_gmc}
Suppose $B=1$, $G\ge 2$, and there are $M$ rollout workers,
each generating one group at a time. Assumptions:
\begin{enumerate}
    \item Every dispatched group is eventually consumed and exactly once.
    After each learner update, the corresponding worker immediately
    starts generating a new group under the updated policy.

    \item At dispatch, $X_k\mid\mathcal G_k\sim\mathcal D$ and
    \[
    Y_k^1,\ldots,Y_k^G\mid X_k,\mathcal G_k
    \overset{\mathrm{i.i.d.}}{\sim}
    \pi_{\theta_{d(k)}}(\cdot\mid X_k).
    \]
    Completion times may depend on the generated responses.

    \item Rewards are bounded by $R$, trajectory lengths by $T$,
    and Assumptions 1 and 2 hold, including along the segments
    joining $\theta_{d(k)}$ and $\theta_{r(k)}$.
    The objective $J$ is $L_J$-smooth and bounded above by $J^\star$.

    \item The GMC cap satisfies the finite cap-switching
    assumption above.
\end{enumerate}

Let $\theta_0$ be deterministic and let $\eta>0$ be constant.
Define
\[
K_0=2R\sqrt{2L\log C},\qquad
K=(1+\epsilon)K_0,
\]
\[
H_\epsilon=2R(1+\epsilon)(4TL\log C+L),
\qquad
H_J=RT\sqrt{2L\log C}.
\]
Then, for every $N\ge 1$,
\begin{align*}
\frac1N\sum_{n=1}^N
\mathbb E\|\nabla J(\theta_n)\|_2^2
&\le
\frac{T(J^\star-J(\theta_0))}{\eta N}
+\frac{TMH_JK_0}{N}\\
&\quad+
T\eta KM(L_JK_0+H_JH_\epsilon)
+\frac{TL_J\eta K^2}{2}.
\end{align*}

In particular, taking $\eta=N^{-1/2}$ and fixing $M$,
\[
\frac1N\sum_{n=1}^N
\mathbb E\|\nabla J(\theta_n)\|_2^2
=O(N^{-1/2}).
\]
\end{theorem}


\begin{proof}
\textbf{Step 1: applying the smoothness of $J$.}
By $L_J$-smoothness of $J$,
\begin{align*}
J(\theta_{n+1})
&\ge J(\theta_n)
+\langle\nabla J(\theta_n),
\theta_{n+1}-\theta_n\rangle
-\frac{L_J}{2}\|\theta_{n+1}-\theta_n\|_2^2\\
&=J(\theta_n)
+\eta\langle\nabla J(\theta_n),g_n\rangle
-\frac{L_J\eta^2}{2}\|g_n\|_2^2\\
&\ge J(\theta_n)
+\eta\langle\nabla J(\theta_n),g_n\rangle
-\frac{L_J\eta^2K^2}{2}.
\end{align*}

Summing over $n=0,\ldots,N-1$, and recalling that
$g_{r(k)}=U_k$, gives
\begin{equation}
J(\theta_N)-J(\theta_0)
\ge
\eta\sum_{k\in\mathcal C_N}
\langle\nabla J(\theta_{r(k)}),U_k\rangle
-\frac{L_J\eta^2NK^2}{2}.
\label{eq:async_descent}
\end{equation}

\textbf{Step 2: decomposing the inner product.}
Recall that $U_k=V_k+e_k$. Adding and subtracting
$\nabla J(\theta_{d(k)})$, we obtain
\begin{align*}
\left\langle\nabla J\left(\theta_{r(k)}\right), U_k\right\rangle & \\
= & \underbrace{\left\langle\nabla J\left(\theta_{d(k)}\right), V_k\right\rangle}_{\text {dispatch-time term }} \\
& +\underbrace{\left\langle\nabla J\left(\theta_{r(k)}\right)-\nabla J\left(\theta_{d(k)}\right), V_k\right\rangle}_{\text {staleness error }} \\
& +\underbrace{\left\langle\nabla J\left(\theta_{r(k)}\right), e_k\right\rangle}_{\text {GMCerror }} .
\end{align*}


\textbf{Step 3: bounding the staleness and GMC errors.}
First, recall that $\|V_k\|_2\le K_0$, as shown
in the virtual-iterate lemma, where
\[
K_0=2R\sqrt{2L\log C}.
\]
By Cauchy--Schwarz, $L_J$-smoothness of $J$
(Lemma 4 in our previous analysis), and
\eqref{eq:parameter_drift},
\begin{align*}
&|\langle
\nabla J(\theta_{r(k)})-\nabla J(\theta_{d(k)}),
V_k\rangle|\\
&\le
\|\nabla J(\theta_{r(k)})
-\nabla J(\theta_{d(k)})\|_2\|V_k\|_2\\
&\le L_JK_0
\|\theta_{r(k)}-\theta_{d(k)}\|_2\\
&\le L_J\eta KK_0\tau_k.
\end{align*}


Summing the staleness error over $k\in\mathcal C_N$
and applying Lemma~\ref{lem:cumulative_delay},
\begin{align*}
&\left|
\sum_{k\in\mathcal C_N}
\langle
\nabla J(\theta_{r(k)})-\nabla J(\theta_{d(k)}),
V_k\rangle
\right|\\
&\le L_J\eta KK_0
\sum_{k\in\mathcal C_N}\tau_k\\
&\le L_J\eta KK_0MN.
\end{align*}


Next, by Lemma 2 (the trajectory score bound)
and the bounded reward assumption,
\[
\|\nabla J(\theta)\|_2
\le RT\sqrt{2L\log C}
=:H_J.
\]

Using Cauchy--Schwarz and the corollary to
Lemma~\ref{lem:cumulative_delay},
\begin{align*}
&\left|
\sum_{k\in\mathcal C_N}
\langle\nabla J(\theta_{r(k)}),e_k\rangle
\right|\\
&\le
\sum_{k\in\mathcal C_N}
\|\nabla J(\theta_{r(k)})\|_2\|e_k\|_2\\
&\le H_J\sum_{k\in\mathcal C_N}\|e_k\|_2\\
&\le H_JH_\epsilon\eta KMN.
\end{align*}

Summing the decomposition from Step 2 over
$k\in\mathcal C_N$, and using the two bounds above
together with Lemma~\ref{lem:cumulative_delay},
we obtain
\begin{align}
&\sum_{k\in\mathcal C_N}
\langle\nabla J(\theta_{r(k)}),U_k\rangle
\nonumber\\
&\ge
\sum_{k\in\mathcal C_N}
\langle\nabla J(\theta_{d(k)}),V_k\rangle
\nonumber\\
&\quad-\eta KMN(L_JK_0+H_JH_\epsilon).
\label{eq:async_error_bound}
\end{align}

\textbf{Step 4. Relating consumed and dispatched groups.}
Since every dispatched group is either consumed or
still outstanding,
\[
\mathcal D_N=\mathcal C_N\sqcup\mathcal P_N.
\]
Hence
\begin{align*}
&\sum_{k\in\mathcal C_N}
\langle\nabla J(\theta_{d(k)}),V_k\rangle\\
&=
\sum_{k\in\mathcal D_N}
\langle\nabla J(\theta_{d(k)}),V_k\rangle\\
&\quad-
\sum_{k\in\mathcal P_N}
\langle\nabla J(\theta_{d(k)}),V_k\rangle.
\end{align*}

Recall that $\|V_k\|_2\le K_0$ from \eqref{eq:gentimeestbound} and
$\|\nabla J(\theta)\|_2\le H_J$ from Step 3.
Since $|\mathcal P_N|=M$, Cauchy--Schwarz gives
\begin{align*}
&\left|
\sum_{k\in\mathcal P_N}
\langle\nabla J(\theta_{d(k)}),V_k\rangle
\right|\\
&\le
\sum_{k\in\mathcal P_N}
\|\nabla J(\theta_{d(k)})\|_2\|V_k\|_2\\
&\le MH_JK_0.
\end{align*}

Therefore,
\begin{align*}
&\sum_{k\in\mathcal C_N}
\langle\nabla J(\theta_{d(k)}),V_k\rangle\\
&\ge
\sum_{k\in\mathcal D_N}
\langle\nabla J(\theta_{d(k)}),V_k\rangle
-MH_JK_0.
\end{align*}

We now take expectations. Recall the dispatch-time
unbiasedness property,
\[
\mathbb E[V_k\mid\mathcal G_k]
=\frac1T\nabla J(\theta_{d(k)}).
\]
Since $\theta_{d(k)}$ is measurable with respect
to $\mathcal G_k$, the tower property gives
\begin{align*}
&\mathbb E[
\langle\nabla J(\theta_{d(k)}),V_k\rangle]\\
&=
\mathbb E\left[
\left\langle
\nabla J(\theta_{d(k)}),
\mathbb E[V_k\mid\mathcal G_k]
\right\rangle
\right]\\
&=\frac1T
\mathbb E\|\nabla J(\theta_{d(k)})\|_2^2.
\end{align*}

There are $M$ initial groups dispatched at $\theta_0$,
and one new group dispatched at 
$\theta_1,\ldots,\theta_N$ each. Thus,
\begin{align*}
&\mathbb E\left[
\sum_{k\in\mathcal D_N}
\langle\nabla J(\theta_{d(k)}),V_k\rangle
\right]\\
&=
\frac MT\mathbb E\|\nabla J(\theta_0)\|_2^2
+\frac1T\sum_{n=1}^N
\mathbb E\|\nabla J(\theta_n)\|_2^2.
\end{align*}

Combining the previous inequalities, we obtain
\begin{align}
&\mathbb E\left[
\sum_{k\in\mathcal C_N}
\langle\nabla J(\theta_{d(k)}),V_k\rangle
\right]\nonumber\\
&\ge
\frac MT\mathbb E\|\nabla J(\theta_0)\|_2^2
+\frac1T\sum_{n=1}^N
\mathbb E\|\nabla J(\theta_n)\|_2^2
-MH_JK_0.
\label{eq:dispatch_bound}
\end{align}


\textbf{Step 5. putting it together.}
Combining \eqref{eq:async_descent} and
\eqref{eq:async_error_bound}, and taking expectations,
we obtain
\begin{align*}
&\mathbb E[J(\theta_N)]-J(\theta_0)\\
&\ge
\eta\mathbb E\left[
\sum_{k\in\mathcal C_N}
\langle\nabla J(\theta_{d(k)}),V_k\rangle
\right]\\
&\quad-\eta^2 KMN(L_JK_0+H_JH_\epsilon)
-\frac{L_J\eta^2NK^2}{2}.
\end{align*}

Now applying \eqref{eq:dispatch_bound},
\begin{align*}
&\mathbb E[J(\theta_N)]-J(\theta_0)\\
&\ge
\frac{\eta M}{T}\mathbb E\|\nabla J(\theta_0)\|_2^2\\
&\quad+
\frac{\eta}{T}\sum_{n=1}^N
\mathbb E\|\nabla J(\theta_n)\|_2^2\\
&\quad-\eta MH_JK_0\\
&\quad-\eta^2 KMN(L_JK_0+H_JH_\epsilon)
-\frac{L_J\eta^2NK^2}{2}.
\end{align*}

Since $J(\theta_N)\le J^\star$, rearranging and
dropping the nonnegative initial-gradient term gives
\begin{align*}
&\frac{\eta}{T}\sum_{n=1}^N
\mathbb E\|\nabla J(\theta_n)\|_2^2\\
&\le J^\star-J(\theta_0)
+\eta MH_JK_0\\
&\quad+\eta^2 KMN(L_JK_0+H_JH_\epsilon)
+\frac{L_J\eta^2NK^2}{2}.
\end{align*}

Dividing by $\eta N/T$, we obtain
\begin{align}
&\frac1N\sum_{n=1}^N
\mathbb E\|\nabla J(\theta_n)\|_2^2
\nonumber\\
&\le
\frac{T(J^\star-J(\theta_0))}{\eta N}
+\frac{TMH_JK_0}{N}
\nonumber\\
&\quad+
T\eta KM(L_JK_0+H_JH_\epsilon)
+\frac{TL_J\eta K^2}{2}.
\label{eq:async_convergence_rate}
\end{align}

Taking $\eta=N^{-1/2}$, for fixed $M$,
\[
\frac1N\sum_{n=1}^N
\mathbb E\|\nabla J(\theta_n)\|_2^2
=O(N^{-1/2}).
\]
\end{proof}